\documentclass[11pt]{article}
\usepackage{mathblatt}
% Präambel des Setzers (werkzeuge/setzer-praeambel.tex, wortgleich); dazu \gteil (Muster T2U)
\makeatletter
\setlength{\mbpfbe}{0mm}% keine Punkte (pruefung.md § 5: nur beim Prüfstein)
\newlength{\gEin}\setlength{\gEin}{0pt}
\let\gSchrift\relax
\renewcommand{\pfkreuzzeile}[1]{\par\noindent\hangindent3.2em\hangafter1\hspace*{1.6em}$\square$\ #1\par}
\newcommand{\gkreuz}[1]{$\square$\ #1\hspace{2em}}
\newcommand{\gkreuzzeile}[2]{\par\noindent\hangindent2.2em\hangafter1\mbox{#1\,$\square$}\ #2\par}% Buchstabe, Option: eine je Zeile (neben Skizze)
% links, was man ansieht, rechts, was man tut; Spaltengrenze überall an derselben Stelle (gemeinsam.md „Satz“)
\newsavebox{\gbox}\newlength{\gLinks}
% v0.10 (Lehrer 13.10.): keine starre Halbierung – die linke Spalte ist so breit wie die Figur (+ 6 mm), die
% Antwort steht direkt daneben; ist die Figur breiter als 0,62 Zeile, steht die Antwort darunter
\newcommand{\gzwei}[2]{\par\smallskip\noindent\sbox{\gbox}{#1}%
  \setlength{\gLinks}{\dimexpr\wd\gbox+6mm\relax}%
  \ifdim\gLinks>0.62\linewidth
    \usebox{\gbox}\par\smallskip\noindent\begin{minipage}{\linewidth}#2\end{minipage}%
  \else
    \begin{minipage}[t]{\gLinks}\vspace{0pt}\usebox{\gbox}\end{minipage}%
    \begin{minipage}[t]{\dimexpr\linewidth-\gLinks\relax}\vspace{0pt}\raggedright #2\end{minipage}%
  \fi\par}
% v0.9 (12.10., pruefung.md § 5 „Marke mit B, Bedingungen rechts“): drittes Argument von pfaufg = Bedingungen
% der Prüfung (v0.10: \pfohnetr = TR durchgestrichen, „FOR“), klein grau im rechten Rand auf der Grundlinie der Nummer, Zeilen mit \\
\newcommand{\gbedrand}[2]{\ifblank{#2}{}{\rlap{\hspace*{\dimexpr#1+2mm\relax}%
  \parbox[t]{15mm}{\raggedright\scriptsize\color{mbgrau}\strut#2}}}}
% Teil eines Bündels (a, b …): Marke links im Rand, Bedingungen rechts im Rand, auf der Zeile des Teils
\newcommand{\gteilrand}[2]{\llap{\mbpfrandmarke{#1}\hspace*{\mbpfnr}}\gbedrand{\linewidth}{#2}}
% Teil eines Bündels: Buchstabe in eigener Spalte, Inhalt eingerückt (\gEin = Einzug, damit die Spaltengrenze von
% \gzwei an derselben Stelle bleibt)
\newenvironment{gbteil}[3]{\par\smallskip\noindent\gteilrand{#2}{#3}\makebox[\mbpfteil][l]{\textbf{#1}}%
  \begin{minipage}[t]{\dimexpr\linewidth-\mbpfteil\relax}\setlength{\gEin}{\mbpfteil}\raggedright\strut\ignorespaces}%
 {\end{minipage}\par}
% Grafik höchstens so breit wie die Spalte (Bündel nebeneinander, Koordinatensystem ohne Karo)
\newsavebox{\gpbox}
\newcommand{\gpassend}[1]{\sbox{\gpbox}{#1}\ifdim\wd\gpbox>\linewidth\typeout{GPASSEND: verkleinert \the\wd\gpbox\space auf \the\linewidth}\resizebox{\linewidth}{!}{\usebox{\gpbox}}\else\usebox{\gpbox}\fi}
% Skizzenfeld (v0.10, Lehrer 13.10.): Karo ab dem linken Rand über die Zeile, „Skizze“ klein grau oben links im Feld
% (wie \pfskizzenfeld der Vorlage); der rechte Rand bleibt den Bedingungen
\newcommand{\gskizze}[1]{\par\smallskip\noindent\begin{tikzpicture}
  \clip (0,0) rectangle ({\linewidth-1mm},{#1*10mm});
  \draw[black!16,line width=0.3pt,step=5mm] (0,0) grid ({\linewidth},{#1*10mm});
  \node[anchor=north west,inner sep=2pt,font=\scriptsize,text=mbgrau,fill=white] at (0,{#1*10mm}) {Skizze};
  \end{tikzpicture}\par}
\RenewDocumentEnvironment{pfaufg}{m m m}{%
  \begin{lrbox}{\mbaufgabebox}\begin{minipage}[t]{\linewidth}%
  \gSchrift\noindent\hspace*{\gEin}\mbpfrandmarke{#1}\gbedrand{\linewidth-\gEin}{#3}%
  \begin{minipage}[t]{\mbpfnr}\strut\textbf{#2}\end{minipage}%
  \begin{minipage}[t]{\dimexpr\linewidth-\gEin-\mbpfnr\relax}\raggedright\strut\ignorespaces}%
 {\end{minipage}%
  \end{minipage}\end{lrbox}%
  \setlength{\mbaufgabehoehe}{\dimexpr\ht\mbaufgabebox+\dp\mbaufgabebox\relax}%
  \ifdefempty{\gkopfnach}{\par\addvspace{5pt}\Needspace*{\mbaufgabehoehe}}{%
    % Stufenkopf, der mit dieser (ersten) Aufgabe zusammenbleibt: Platz für beide verlangen, dann den Kopf setzen
    % (12.10., Quadratische: ein Kopf stand allein unten, die Aufgabe mit Koordinatensystem auf der nächsten Seite)
    \par\Needspace*{\dimexpr\mbaufgabehoehe+40pt\relax}\gkopfnach\gdef\gkopfnach{}\par\addvspace{5pt}}%
  \noindent\usebox{\mbaufgabebox}\par\addvspace{5pt}}
% Stufenkopf vormerken: \gkopfvor{<Kopf>} – die nächste pfaufg setzt ihn, sobald beide zusammen Platz haben
\def\gkopfnach{}
\newcommand{\gkopfvor}[1]{\gdef\gkopfnach{#1}}
% Teilaufgabe mit Skizze links in gzwei: Buchstabe in eigener Spalte, darunter Text und Skizze
\newcommand{\gteil}[2]{\makebox[\mbpfteil][l]{\textbf{#1}}\parbox[t]{\dimexpr0.5\textwidth-\gEin-\mbpfnr-5mm-\mbpfteil\relax}{\raggedright #2}}
% Stufenkopf der Vorlage (mathblatt.sty 2026-10-12: Linie, Abstand, 13 pt, \pfjahre) ohne ihren Vorhalt von
% 14 Zeilen: der Setzer fordert vor dem Kopf 0,16 Seitenhöhe (Kopf + erste Aufgabe), mehr nicht
\setlength{\mbpfstufevorhalt}{0pt}
\makeatother
\newcommand{\karo}[1]{\par\smallskip\noindent\begin{tikzpicture}
  \clip (0,0) rectangle ({\linewidth-1mm},{#1*10mm});
  \draw[black!16,line width=0.3pt,step=5mm] (0,0) grid ({\linewidth},{#1*10mm});
  \end{tikzpicture}\par}
\newcommand{\kk}[1]{$\square$\,#1\hspace{0.9em}}
\newcommand{\linie}[1]{\,{\color{black!45}\rule[-1pt]{#1}{0.4pt}}\,}
\newcommand{\pkt}{\textperiodcentered{}}
% v0.7 (12.10.): Schülerfrage unter dem Stufenkopf; Ankreuzoption als Bild (Buchstabe und Kästchen darunter)
\newcommand{\gfrage}[1]{\par\nobreak\noindent\hspace*{\mbpfnr}{\small\itshape #1}\par\nobreak\addvspace{2pt}}
\newlength{\gbildbreite}\setlength{\gbildbreite}{0.32\linewidth}% Setzer setzt sie je Reihe (vier Bilder: 0,24)
\newcommand{\gbildoption}[3]{\begin{minipage}[t][#3][b]{\gbildbreite}\centering #2\par\smallskip\mbox{#1\,$\square$}\end{minipage}}
% Lineare (12.10.): Bild mit Buchstabe ohne Kästchen (Diagramm zuordnen); Aussage mit „wahr / falsch“ zum Ankreuzen
\newcommand{\gbildname}[3]{\begin{minipage}[t][#3][b]{\gbildbreite}\centering #2\par\smallskip\mbox{\textbf{#1}}\end{minipage}}
\newcommand{\gwahrfalsch}[2]{\par\noindent\hangindent2.2em\hangafter1\makebox[2.2em][l]{(#1)}#2\hfill\mbox{wahr $\square$\quad falsch $\square$}\par}
\begin{document}

\pfheftstil
\pflaufkopf[P10 \pkt{} Prüfung]{Parabel: Scheitelpunkt und Skizze}{XU62}\pfseitevon
\pfheftkopf{Parabel: Scheitelpunkt und Skizze}{}
\pfstufenzeile{Wertetabelle oder Skizze \pkt{} Scheitelpunkt angeben \pkt{} Punkt auf der Parabel prüfen \pkt{} Gleichung in Scheitelpunktform \pkt{} Lage zweier Parabeln begründen}

% ===================== Stufe: Wertetabelle oder Skizze =====================
\gkopfvor{\pfstufekopf{Wertetabelle oder Skizze}{\pfjahre{26 \pkt{} 24}{15}}}
% kein Einstieg: keine Schülerfrage in der Gliederung (Schülerfrage, -Fälle, -Form, -Weg, -Skizze; oder Frage) – Einstieg nur von dort, nie aus der Bank (pruefung.md § 5, 12.10.)
% 2026-FOR-B1e-m1 (Original 2026-FOR-B1e, Schwere I, rang 1; herausgelöste Fassung: keine)
\begin{pfaufg}{’26~\pkt{}~1e}{1.}{Basis}
Kreuze die Wertetabelle an, die zu $y = 2x^2$ gehört.
\par\smallskip\noindent\setlength{\gbildbreite}{\dimexpr0.96\linewidth/3\relax}\sbox{\gbox}{\vbox{{\setlength{\mbzell}{5mm}\setlength{\tabcolsep}{4pt}\wertetabelle[8,2,0,2,8]{x}{y}{-2,-1,0,1,2}}\par{\setlength{\mbzell}{5mm}\setlength{\tabcolsep}{4pt}\wertetabelle[-4,-2,0,2,4]{x}{y}{-2,-1,0,1,2}}\par{\setlength{\mbzell}{5mm}\setlength{\tabcolsep}{4pt}\wertetabelle[16,4,0,4,16]{x}{y}{-2,-1,0,1,2}}}}\gbildoption{A}{{\setlength{\mbzell}{5mm}\setlength{\tabcolsep}{4pt}\wertetabelle[8,2,0,2,8]{x}{y}{-2,-1,0,1,2}}}{\dimexpr\ht\gbox/3+\baselineskip\relax}\hfill\gbildoption{B}{{\setlength{\mbzell}{5mm}\setlength{\tabcolsep}{4pt}\wertetabelle[-4,-2,0,2,4]{x}{y}{-2,-1,0,1,2}}}{\dimexpr\ht\gbox/3+\baselineskip\relax}\hfill\gbildoption{C}{{\setlength{\mbzell}{5mm}\setlength{\tabcolsep}{4pt}\wertetabelle[16,4,0,4,16]{x}{y}{-2,-1,0,1,2}}}{\dimexpr\ht\gbox/3+\baselineskip\relax}\par
\fusshilfe{\mbox{1:~A}}
\end{pfaufg}
% 2024-OS-K3b-m1 (Original 2024-OS-K3b, Schwere I, rang 4; herausgelöste Fassung: keine)
\begin{pfaufg}{nach~P10\\’24~\pkt{}~3b}{2.}{}
Die Gerade $f$ ist der Graph von $f(x) = 2x + 3$. Die Parabel $p$ ist der Graph von $p(x) = x^2 - 5$. Skizziere die Parabel $p$ in das Koordinatensystem. Gib ihren Scheitelpunkt an: $S(\ \ |\ \ )$.
\par\smallskip\noindent\gpassend{\begin{ksys}[xmin=-4,xmax=4,ymin=-6,ymax=6,karo=0.5]\end{ksys}}\par
\fusshilfe{\mbox{2:~S(0|$-$5)}}
\end{pfaufg}

% ===================== Stufe: Scheitelpunkt ablesen =====================
\gkopfvor{\pfstufekopf{Scheitelpunkt angeben}{\pfjahre{26 \pkt{} 22}{21 \pkt{} 20 \pkt{} 18 \pkt{} 17}}}
% kein Einstieg: keine Schülerfrage in der Gliederung (Schülerfrage, -Fälle, -Form, -Weg, -Skizze; oder Frage) – Einstieg nur von dort, nie aus der Bank (pruefung.md § 5, 12.10.)
% 2026-FOR-K5c-m1 (Original 2026-FOR-K5c, Schwere I, rang 2; herausgelöste Fassung: keine)
\begin{pfaufg}{nach~P10\\’26~\pkt{}~5c}{3.}{}
Die verschobene Normalparabel $p$ ist der Graph von $p(x) = (x-3)^2 - 1$ und ist im Koordinatensystem gezeichnet. Gib den Scheitelpunkt von $p$ an: $S(\ \ |\ \ )$.
\gzwei{\begin{ksys}[xmin=-1,xmax=6,ymin=-2,ymax=6,karo=0.5]\parabel{1}{3}{-1}{p}\end{ksys}}{\karo{3}}
\fusshilfe{\mbox{3:~S(3|$-$1)}}
\end{pfaufg}
% 2022-OS-K3b-m1 (Original 2022-OS-K3b, Schwere I, rang 5; herausgelöste Fassung: keine)
\begin{pfaufg}{nach~P10\\’22~\pkt{}~3b}{4.}{}
Die lineare Funktion $f$ hat die Gleichung $f(x) = -2x + 5$. Die verschobene Normalparabel $p$ hat die Gleichung $p(x) = (x-4)^2 - 6$. Gib den Scheitelpunkt $S$ von $p$ an.
\hfill\linie{30mm}
\fusshilfe{\mbox{4:~S(4|$-$6)}}
\end{pfaufg}

% ===================== Stufe: Punkt auf der Parabel prüfen =====================
% Formel-Frage (einmal je Blatt, erste Stufe mit Formel; pruefung.md § 5)
\gkopfvor{\pfstufekopf{Punkt auf der Parabel prüfen}{\pfjahre{}{20 \pkt{} 18 \pkt{} 14}}\gfrage{Wo steht die Formel in deiner Formelsammlung? Seite \linie{12mm}}}
% kein Einstieg: keine Schülerfrage in der Gliederung (Schülerfrage, -Fälle, -Form, -Weg, -Skizze; oder Frage) – Einstieg nur von dort, nie aus der Bank (pruefung.md § 5, 12.10.)
% Teile eines Originals (2020-OS-K3): 2020-OS-K3b-m1, 2020-OS-K3c-m1
\begin{pfaufg}{}{5.}{}
Die Parabel zu $f$ mit $y = (x+3)^2 - 5$ ist im Koordinatensystem gezeichnet.
% 2020-OS-K3b-m1 (Original 2020-OS-K3b, Schwere I, rang 10; herausgelöste Fassung: keine)
\begin{gbteil}{a)}{\raisebox{\dimexpr-\height+\ht\strutbox\relax}{\shortstack[r]{nach~P10\\’20~\pkt{}~3b}}\hspace*{\mbpfnr}}{}
Genau einer dieser Punkte liegt nicht auf dem Graphen von $f$. Kreuze ihn an:
\gzwei{\begin{ksys}[xmin=-7,xmax=2,ymin=-6,ymax=5,karo=0.5]\parabel{1}{-3}{-5}{f}\end{ksys}}{\vspace{2mm}\pfkreuzab{A}{$A(-5\,|\,-1)$}\pfkreuzab{B}{$B(-1\,|\,-1)$}\pfkreuzab{C}{$C(0\,|\,4)$}\pfkreuzab{D}{$D(-2{,}5\,|\,-4{,}5)$}}
\end{gbteil}
% 2020-OS-K3c-m1 (Original 2020-OS-K3c, Schwere I, rang 11; herausgelöste Fassung: keine)
\begin{gbteil}{b)}{\raisebox{\dimexpr-\height+\ht\strutbox\relax}{\shortstack[r]{nach~P10\\’20~\pkt{}~3c}}\hspace*{\mbpfnr}}{}
Der Punkt $A(2\,|\,y)$ liegt auf dem Graphen von $f$ mit $y = (x+3)^2 - 5$. Berechne $y$.
\karo{2}
\par\smallskip\par\noindent \hfill y: \linie{30mm}\hphantom{\,cm²}
\end{gbteil}
\fusshilfe{\mbox{5:~a)~D; b)~20}}
\end{pfaufg}
% 2018-OS-K5a-m1 (Original 2018-OS-K5a, Schwere II, rang 17; herausgelöste Fassung: keine)
\begin{pfaufg}{nach~P10\\’18~\pkt{}~5a}{6.}{}
Im Koordinatensystem ist die Parabel $p$ mit $p(x) = x^2 + 2x - 2$ gezeichnet. Entscheide für jede Aussage, ob sie wahr oder falsch ist:
\gwahrfalsch{1}{Der Punkt $(-3\,|\,1)$ liegt auf $p$.}\gwahrfalsch{2}{$p$ ist die Normalparabel, um 1 nach links und um 3 nach unten verschoben.}\gwahrfalsch{3}{$p$ schneidet die y-Achse im Punkt $(-2\,|\,0)$.}
\par\smallskip\noindent\begin{ksys}[xmin=-5,xmax=4,ymin=-4,ymax=7,karo=0.5]\parabel{1}{-1}{-3}{p}\end{ksys}\par
\fusshilfe{\mbox{6:~1; Aussage 2 wahr, Aussage 3 falsch (y-Achse bei (0|$-$2))}}
\end{pfaufg}
% 2014-OS-K7a-m1 (Original 2014-OS-K7a, Schwere II, rang 21; herausgelöste Fassung: keine)
\begin{pfaufg}{nach~P10\\’14~\pkt{}~7a}{7.}{FOR}
Im Koordinatensystem sind die Parabel $p$ mit $p(x) = -x^2$ und die Gerade $g$ mit $g(x) = 3x - 4$ gezeichnet. Untersuche mit einer Rechnung, ob $S(-4\,|\,-16)$ ein gemeinsamer Punkt von $p$ und $g$ ist.
\gzwei{\begin{ksys}[xmin=-6,xmax=6,ymin=-6,ymax=6,karo=0.5]\parabel{-1}{0}{0}{p}\gerade[2.5]{3}{-4}{g}\end{ksys}}{\karo{3}}
\fusshilfe{\mbox{7:~$-$16; $-$16}}
\end{pfaufg}

% ===================== Stufe: Scheitelpunktform angeben =====================
\gkopfvor{\pfstufekopf{Gleichung in Scheitelpunktform}{\pfjahre{25 \pkt{} 23}{20 \pkt{} 18 \pkt{} 17 \pkt{} 16 \pkt{} 15}}}
% kein Einstieg: keine Schülerfrage in der Gliederung (Schülerfrage, -Fälle, -Form, -Weg, -Skizze; oder Frage) – Einstieg nur von dort, nie aus der Bank (pruefung.md § 5, 12.10.)
% 2025-OS-K5b-m1 (Original 2025-OS-K5b, Schwere I, rang 3; herausgelöste Fassung: keine)
\begin{pfaufg}{nach~P10\\’25~\pkt{}~5b}{8.}{}
Die quadratische Funktion $p$ hat die Gleichung $p(x) = x^2 - 8x + 13$. Ihr Graph ist im Koordinatensystem gezeichnet. Gib den Scheitelpunkt $S$ von $p$ an: $S(\ \ |\ \ )$. Schreib die Gleichung von $p$ in Scheitelpunktform: $p(x) = \ldots$
\gzwei{\begin{ksys}[xmin=-1,xmax=8,ymin=-4,ymax=6,karo=0.5]\parabel{1}{4}{-3}{p}\end{ksys}}{\karo{4}}
\fusshilfe{\mbox{8:~S(4|$-$3); (x $-$ 4)² $-$ 3}}
\end{pfaufg}
% 2023-OS-K4b-m1 (Original 2023-OS-K4b, Schwere II, rang 6; herausgelöste Fassung: keine)
\begin{pfaufg}{nach~P10\\’23~\pkt{}~4b}{9.}{FOR}
Gib den Scheitelpunkt der gezeichneten Normalparabel $p$ an: $S(\ \ |\ \ )$. Schreib eine Gleichung von $p$ in Scheitelpunktform.
\gzwei{\begin{ksys}[xmin=-5,xmax=5,ymin=-4,ymax=6,karo=0.5]\parabel{-1}{1}{5}{p}\end{ksys}}{\karo{4}}
\fusshilfe{\mbox{9:~S(1|5); $-$(x $-$ 1)² + 5}}
\end{pfaufg}

% ===================== Stufe: Lage zweier Parabeln ohne Rechnung begründen =====================
\gkopfvor{\pfstufekopf{Lage zweier Parabeln begründen}{\pfjahre{26}{14}}}
% kein Einstieg: keine Schülerfrage in der Gliederung (Schülerfrage, -Fälle, -Form, -Weg, -Skizze; oder Frage) – Einstieg nur von dort, nie aus der Bank (pruefung.md § 5, 12.10.)
% 2026-FOR-K5d-m1 (Original 2026-FOR-K5d, Schwere III, rang 7; herausgelöste Fassung: keine)
\begin{pfaufg}{nach~P10\\’26~\pkt{}~5d}{10.}{FOR}
Die verschobene Normalparabel $p$ hat die Gleichung $p(x) = (x-3)^2 - 1$. Eine zweite Parabel $q$ hat die Gleichung $q(x) = -(x-3)^2 - 1$. Gib an, wie viele Punkte $p$ und $q$ gemeinsam haben. Begründe ohne zu rechnen.
\karo{3}
\fusshilfe{\mbox{10:~genau 1 gemeinsamer Punkt}}
\end{pfaufg}

\pfblattende{11}{26 FOR & 25 & 24 & 23 & 22 & 21 & 20 & 19 & 18 & 17 & 16}{1e 5c 5d & 5b & 3b & 4b & 3b & 1d & 3a 3b 3c 3d &  & 5a 5b 5c 5d & 5c 5e & 1g}%
{\pfausschnitt{}{Parabel: Scheitelpunkt und Skizze}{Quadratische Gleichung lösen}}
\end{document}